\section{Introduction}
\label{sec:introduction}

Media bias, the systematic slanting of news coverage through framing, omission, or word choice, poses a significant threat to informed public discourse \cite{hamborg2019automated}.
As online news consumption grows, automated methods for detecting biased reporting have become a priority in the NLP community \cite{rodrigo2024systematic}.
Most existing approaches analyze article content directly, relying on lexical, syntactic, or semantic features to classify bias at the document or sentence level.
While effective, these methods face two fundamental limitations: they require expensive annotation efforts, and they capture only the producer side of bias, ignoring how audiences perceive and react to biased content.

Social news platforms, where users collectively curate, discuss, and evaluate news articles, offer a complementary perspective.
On these platforms, user interactions (votes, comments, karma scores) encode collective judgments about content quality, credibility, and controversiality.
If biased articles systematically elicit different interaction patterns than unbiased ones, these behavioral signals could serve as scalable proxies for perceived bias, reducing dependence on manual annotation.

In this work, we investigate this hypothesis using Men\'eame (\texttt{meneame.net}), Spain's largest social news aggregator, analogous to Reddit or Digg but focused on the Spanish-language media ecosystem.
Men\'eame provides a rich setting for this study: it aggregates content from thousands of outlets, supports community voting and threaded comments with karma scores, and has operated continuously since 2005, offering a longitudinal view of media consumption patterns.

We make the following contributions:
\begin{enumerate}
    \item We construct and release a large-scale dataset of 14,995 news articles with automatic bias labels and 38 interaction features derived from 13.2 million comments by 96,034 users across 2,868 media outlets (2005--2021).
    \item We conduct a comprehensive statistical analysis demonstrating that biased articles elicit significantly different user engagement patterns in karma distributions, comment sentiment, reaction speed, and thread dynamics.
    \item We perform community analysis on the user-outlet bipartite graph, revealing two distinct media ecosystems with different bias exposure profiles.
    \item We evaluate the predictive power of interaction-only features for bias detection, establishing baselines for interaction-aware bias classification.
\end{enumerate}

The remainder of the paper is organized as follows.
Section~\ref{sec:related_work} reviews related work on media bias detection and user interaction analysis.
Section~\ref{sec:dataset} describes the dataset construction pipeline.
Section~\ref{sec:methods} details our analysis methodology.
Section~\ref{sec:results} presents experimental results.
Section~\ref{sec:discussion} discusses findings, limitations, and implications.
Section~\ref{sec:conclusion} concludes with future directions.
